\chapter{Introduction and Project Vision}

\section{Project Overview}

Unify is a locally deployed system that checks whether a proposed code change
actually does what the ticket and design it claims to implement say it should,
and shows the code and tests that support that conclusion.

A software team's record of a single piece of work is normally scattered across
three tools. The intended appearance and behaviour live in a design tool. The
requirement and its workflow state live in an issue tracker. The implementation
lives in a repository. Each tool versions its own record independently and none
of them understands what a change in one implies for the others. Unify
subscribes to all three, resolves which records refer to the same unit of work,
stores them with their versions and their provenance, and evaluates each
proposed change against the intent it is linked to.

The output is not a score and not a recommendation to merge. For each acceptance
criterion, the system reports whether the available evidence supports it,
contradicts it, supports it only in part, or is insufficient to decide, and every
decisive finding is accompanied by the file, line range and artifact version it
was derived from. A finding that cannot cite its evidence is not shown as a
finding.

All model inference runs on infrastructure the team owns. No repository content,
ticket text or design data is transmitted to a third-party model interface.

\section{Problem Statement}

Work drifts out of alignment across the tools that hold it, and the
inconsistency is discovered late.

A designer updates a screen after the corresponding implementation has already
been written. A ticket is marked complete while the code that implements it is
still changing. A pull request is reviewed against acceptance criteria that were
edited after the branch was cut. In each case the individual tools are behaving
correctly: each has an accurate record of its own artifact. What is missing is
any representation of the relationship between those artifacts, and therefore
any mechanism by which a change to one can raise a question about another.

The concrete decision that fails is a code review. A reviewer opening a pull
request must determine whether the change still matches the intent it claims to
implement. Doing so requires opening the tracker for the acceptance criteria,
opening the design tool and checking whether the approved design has moved since
the branch was created, reading the diff, and forming a judgement about whether
the tests actually exercise the criterion in question. That reasoning is
performed from memory across several interfaces, and no artifact records what
evidence supported the decision, so it is lost the moment the change is merged.

The consequence is that mismatches surface during testing, at deployment, or in
a status meeting, after the cost of building the wrong thing has already been
incurred.

\section{Problem Context and Motivation}

The problem is a consequence of tool specialisation. Teams adopt a design tool,
an issue tracker and a version control platform because each is substantially
better at its own task than a single integrated product would be. The cost of
that choice is that the relationships between their contents exist only in the
heads of the people doing the work.

Existing integrations address a narrower problem than the one described here.
Field synchronisation tools copy values between systems but do not reason about
whether the copied values remain consistent with one another. Pairwise
integrations connect two specific tools and notify on change, but see neither the
code nor the tests. Recent commercial code-review assistants do compare a pull
request against a linked issue, which overlaps directly with part of this
project's scope; they operate as hosted services, they do not preserve the
versioned evidence behind a conclusion, and they do not treat uncertainty as a
reportable outcome.

To confirm that the problem is real rather than assumed, the team consulted a
practising senior software engineer about his team's workflow. Three findings
shaped the design. Relationships between pull requests and tickets are currently
tracked manually by technical leads. Implementations built against
subsequently-changed designs have occurred repeatedly and have caused
rework. And the team already receives more automated notifications than it
reads, so a system that raises alerts of uncertain accuracy will be disabled
rather than tuned.

That third finding is the one that most constrains the design. It establishes
that precision matters more than coverage, and it is the reason the system is
built to abstain when evidence is insufficient rather than to guess.

\section{Project Vision}

A reviewer should be able to answer the question ``does this change do what it
says'' from one place, with the evidence attached, and the system should say
plainly when it cannot answer rather than producing a confident guess.

The longer-term intent is that the reasoning behind a review decision stops
being ephemeral. Today it exists only while the reviewer holds it; afterwards
the repository records that a change was approved but not what was checked.
Unify produces a durable, versioned evidence case for each change, which means
the basis for a past decision can be re-read against the artifact versions that
were current when it was made.

\section{Project Objectives}

\begin{enumerate}
  \item Resolve the identity of a single unit of work across a design source, a
        work-item source and a change source, recording which version of each
        artifact a relationship was established against, and achieve high
        precision on links derived from explicit references.
  \item Process notifications from those sources such that duplicate and
        out-of-order delivery converge to the same final state, and an event
        the source never delivered is recovered by reconciliation rather than
        silently missed.
  \item Extract, for a declared set of languages, the changed files, symbols,
        dependency neighbourhood and test evidence associated with a proposed
        change.
  \item Verify each atomic acceptance criterion against a deterministically
        assembled evidence packet, using a bounded local inference loop, and
        return a structured verdict with cited evidence spans.
  \item Prevent any verdict whose citations cannot be resolved against the
        evidence packet from being presented as a decisive finding.
  \item Anchor functional requirements to specific symbols and line ranges,
        pinned to a commit, and maintain those anchors as the code changes,
        reporting the proportion that survive.
  \item Evaluate the system against a manually constructed and independently
        labelled set of cases, and against stated baselines including the
        current manual workflow, reporting performance per class rather than as
        a single aggregate.
\end{enumerate}

\section{Project Scope}

\subsection*{Included}

Read-only ingestion from one design-intent source, one work-item source and one
change source. Deterministic and manually confirmed linking across those
sources. An immutable, versioned artifact store. A typed evidence graph carrying
provenance on every relationship. Deterministic policies covering design-version
drift, workflow-state inconsistency, overlapping change scope and
requirement--implementation divergence. Code and test analysis for the declared
languages. Criterion decomposition and deterministic context assembly. A bounded
local verification loop. A deterministic verdict and provenance gate. A reviewer
interface presenting evidence cases with resolution and audit history.
Role-based routing of findings and outbound delivery to one channel.

\subsection*{Excluded}

Autonomous code generation or repair. Any claim of general program correctness.
Support for languages beyond those declared. Unrestricted conversational
interaction with the system. Two-way write-back of data to the source systems.
Any hosted model interface in the committed system. Pixel-level certification
that an implementation matches a design. Additional source integrations beyond
those bound at the proposal stage.

The exclusions are stated as firmly as the inclusions because several of them
describe capabilities the system could plausibly be extended to offer, and an
unstated boundary tends to be assumed away during evaluation.

\section{Project Modules}

\begin{longtable}{>{\raggedright\arraybackslash}p{0.26\textwidth}>{\raggedright\arraybackslash}p{0.68\textwidth}}
\toprule
Module & Responsibility \\
\midrule
\endhead
Source adapters
  & Authenticated read-only access per source role; notification receipt with
    polling fallback; translation of vendor payloads into the normalised event
    envelope. \\
Event reliability
  & Idempotency by stable event identifier; per-resource checkpointing;
    reconciliation of undelivered events; version guards; stale-source
    detection. \\
Artifact and version store
  & Immutable snapshots of every artifact version, so that evidence cited by a
    finding remains readable after the source changes. \\
Code and test analyser
  & Base and head resolution; changed files and diff hunks; changed and
    affected symbols; bounded dependency neighbourhood; test evidence; explicit
    markers where a relation cannot be resolved statically. \\
Evidence graph
  & Typed nodes and edges with source class, confidence, resolver and validity
    interval; requirement anchors; deterministic and manual linking. \\
Policy and rule engine
  & Versioned declarative policies evaluated incrementally against graph
    deltas; criterion invalidation and queueing. \\
Context builder
  & Deterministic per-criterion evidence packets under a fixed token budget,
    with recorded truncation and a stored manifest. \\
Local verifier
  & Bounded multi-pass inference over one packet, producing a structured verdict
    with cited spans, executed on team-operated infrastructure. \\
Verdict and provenance gate
  & Schema validation, citation resolution, eligibility and threshold checks,
    and downgrade of any decisive output that fails them. \\
Reviewer surface and routing
  & Evidence case presentation; acknowledge, resolve, dispute and suppress
    actions; role-based routing; audit history. \\
\bottomrule
\end{longtable}

\section{Stakeholders and User Classes}

\begin{longtable}{>{\raggedright\arraybackslash}p{0.22\textwidth}>{\raggedright\arraybackslash}p{0.34\textwidth}>{\raggedright\arraybackslash}p{0.38\textwidth}}
\toprule
User class & Characteristics & What they need from the system \\
\midrule
\endhead
Code reviewer or technical lead \emph{(primary)}
  & Technically expert; time-constrained; accountable for the decision to
    approve.
  & To establish quickly whether a change is adequately supported by its linked
    intent, with evidence they can inspect rather than trust. \\
Developer
  & Author of the change; familiar with the code, less so with the design
    history.
  & Notification when the design or requirement they are implementing against
    has changed, and visibility of what evidence their change is missing. \\
Designer
  & Owns the design source; not necessarily a reader of code.
  & To be told when an implementation diverges from an approved structured
    constraint, and to have their design treated as versioned intent. \\
Product owner or manager
  & Owns acceptance criteria and workflow state; non-technical in many teams.
  & Project-level visibility of unresolved evidence without commit-level noise;
    approval rather than authorship of configuration. \\
\bottomrule
\end{longtable}

A secondary consideration applies to the product owner. Reducing that role's
operational involvement is an explicit design objective arising from evaluation
feedback: configuration that would previously have been authored by hand is
derived from the project's own observed behaviour and presented for approval,
and link confirmation is routed to the code owner rather than to the manager.

\section{Success Criteria}

Success is defined by measurement against stated baselines rather than by
completion of features. The criteria below are the targets the system is built
to meet; the figures are those baselined at the proposal stage and are
reconfirmed after the technical spike.

\begin{longtable}{>{\raggedright\arraybackslash}p{0.46\textwidth}>{\raggedright\arraybackslash}p{0.46\textwidth}}
\toprule
Criterion & Target \\
\midrule
\endhead
Precision of links derived from explicit references & At or above 0.95 \\
Retrieval of relevant evidence for a criterion, within the top five candidates
  & Recall at or above 0.90 \\
Precision of findings reporting a contradiction & At or above 0.90 \\
Rate of confident assertions unsupported by valid evidence & At or below 2\% \\
Provenance completeness for every decisive finding & 100\% \\
Correct abstention on cases labelled ambiguous & At or above 80\% \\
Convergence of final state under duplicate and out-of-order delivery & 100\% on
  the controlled event suite \\
Equivalence between incremental and full evaluation & At or above 99\% \\
Agreement between independent annotators on the evaluation set & At or above the
  stated coefficient threshold, or the labelling guide is revised and the set
  relabelled \\
\bottomrule
\end{longtable}

Beyond these component measures, the system is assessed at the workflow level
against the current manual process: whether a reviewer reaches a correct
decision more often, or locates the relevant evidence faster, without an
increase in findings they dismiss.

\section{Constraints and Assumptions}

\subsection*{Constraints}

\textbf{Deployment and privacy.} The system runs within a single self-hosted
deployment boundary, and artifact content does not leave it. This rules out
hosted model interfaces in the committed system and bounds model selection to
what the available hardware can serve.

\textbf{Source system behaviour.} Notification coverage, delivery guarantees and
rate limits are set by the external products and cannot be altered. The design
accommodates them rather than assuming they are uniform.

\textbf{Language coverage.} Code analysis requires a parser, and parsers are
language-specific. The committed analyser supports a declared set of language
families, and extending it is additional work rather than configuration.

\textbf{Evaluation data.} No public dataset links a design artifact, a work item
and a code change for the same unit of work. The evaluation set is therefore
constructed and labelled by the team, which bounds its size and requires that
its construction be independently validated.

\textbf{Team and schedule.} Three members across two semesters, each owning one
technical stream together with the evidence demonstrating it.

\subsection*{Assumptions}

That teams reference tickets in pull requests and link designs to tickets often
enough for deterministic linking to establish a useful proportion of
relationships, with manual linking covering the remainder. That acceptance
criteria, where present, are specific enough to be decomposed into atomic
statements. That a structured design constraint can be confirmed by a designer
or product owner, since unstructured visual similarity is explicitly not treated
as sufficient evidence. And that a reviewer will act on a finding that cites its
evidence, which is the premise the entire output format rests on.

\section{Ethical, Legal and DEI Considerations}

\textbf{Privacy and data minimisation.} The system processes workflow metadata
and the content necessary to verify a criterion. It does not perform
individual-level performance measurement, and the audit record exists to make
findings re-readable rather than to attribute productivity to people. Because
inference is local, source code and design material remain inside the team's own
infrastructure, which also makes the system usable by organisations whose
policies prevent sending code to an external service.

\textbf{Access and credentials.} All access to source systems is read-only and
least-privilege. Credentials are encrypted at rest and a revocation path is
documented. The system never writes to the connected tools.

\textbf{Treatment of untrusted content.} Issue text, commit messages and design
metadata are treated as untrusted input. The verifier does not act on
instructions embedded in artifact content, and model output is constrained to a
schema and validated against a source allowlist before it can influence
anything.

\textbf{Accountability for decisions.} The system does not approve or reject
changes. It presents evidence and the human decides. This is a deliberate
ethical position rather than a limitation of the implementation: the highest
autonomy level, in which a model writes to the system of record or resolves a
finding on its own, is prohibited in this project.

\textbf{Inclusion across roles.} Designers, developers and product owners
currently inspect the same work through tools that are legible to some of them
and opaque to others. Presenting one evidence case, routed to the role
accountable for it, is intended to make cross-disciplinary review possible
without requiring every participant to read code.

\textbf{Legal and licensing.} Integrations operate within the developer terms of
each connected product. The open-weight model, its serving runtime and the
project's dependencies are used in accordance with their licences, which are
recorded as part of model selection rather than assumed. Evaluation repositories
are selected for permissive licensing, and no retained artefact contains
credentials or secrets.
